\section{Conclusion}

We presented AutoShotV2, a compact refinement framework for short-video shot
boundary detection. By freezing the NAS-searched AutoShot backbone and
retraining a calibrated two-branch classification head, the method addresses
the extreme class imbalance characteristic of short videos without repeating
the computationally expensive architecture search. Temperature scaling and
Gaussian smoothing provide the calibrated decision pipeline used at
deployment, reaching F1 scores of \PaperDeployShotFOne{} on SHOT,
\PaperDeployBBCFOne{} on BBC, and \PaperDeployClipFOne{} on ClipShots with a
single fixed configuration.

The controlled replication clarifies where the improvement comes from:
Focal Loss and many-hot supervision leave the BCE control statistically
unchanged, whereas calibrated post-processing yields paired-bootstrap gains
on SHOT and ClipShots; a three-seed study bounds training-seed F1
variability below 0.5\% on every dataset, and class-balanced calibration
error shows why very small aggregate ECE values should not be read as
uniformly calibrated boundary probabilities. Every reported number is
regenerated and checked from tracked artifacts with exact cache identities
and frozen selection manifests.

Cross-domain analysis identifies the main open problem: on open-world
ClipShots the dominant loss is precision, as visual effects and domain shift
create false-positive peaks that survive smoothing. Future work should
reserve an independent calibration split, extend calibration to
boundary-conditional estimates, and explore domain-adaptive components ---
lightweight adapters on the frozen backbone or per-domain post-processing
selection --- together with multimodal cues such as audio for robustness
where visual evidence alone is ambiguous.
